% =============================================================================
% ADE lattice polygons (Alexeev-Thompson)
% =============================================================================
% One object, the toric ADE surface (Y, C) of a given shape:
%   \pic (Q) {ade polygon={<letter>}{<rank>}{<variant>}};
% <letter> is A, D or E; <variant> is one of
%   pure        A_{2n-1}, A^-_{2n-2}, D_{2n}, D^-_{2n-1}, -E_6^-, -E_7, -E_8^-
%               (AT Table 1; the right decoration follows from the rank)
%   left-short  -A^-_{2n-1} and -A_{2n}: A with the left side short (AT Table 1
%               row -A^-_{2n-3}, and its mirror A^-_{2n} of AT Remark 3.9)
%   prime       'A_{2n-1}, 'A^-_{2n-2}, D'_{2n} (AT Table 4)
%   both-prime  'A'_{2n-1} (AT Table 4)
%   affine      \tilde D_{2n} (2n >= 6), \tilde E_7, \tilde E_8^- (AT Table 1)
%
% Source of the data: Alexeev-Thompson, "ADE surfaces and their moduli"
% (arXiv:1712.07932), Table 1 (pure shapes), Table 4 (toric primed shapes),
% Notation 3.7 (the Dynkin diagram on the polygon) and Notations 3.8, 3.21
% (the decorated symbol). The research preamble implements the same polygons
% in dzack_research.preamble.categories.schemes.ade_surfaces
% (ADELogPairs.at21(letter, rank, variant=..., affine=...)); its variants map
% here as pure->pure (A odd, D even, E), short->pure (A even, D odd),
% both-short->left-short (A odd), prime->prime, affine=True->affine.
% Checked against the preamble: polygon, p*, and the lengths of the sides
% through p* agree for every shape both implement, except 'A^-_{2n-2}, where
% the preamble returns the polygon of 'A_{2n-3} (vertex (2n-4,0), not
% (2n-3,0) as in AT Table 4), and A^-_0, which at21 rejects (rank >= 1)
% although AT Table 1 lists it (min(n) = 1).
%
% What the object draws, all derived from the polygon Q and p*:
% - the lattice [0,w]x[0,h] with points named <i>-<j>, and Q filled;
% - the boundary divisor C (the sides through p*) as `long side` (even lattice
%   length) or `short side` (odd), labelled with its lattice length; the other
%   sides C' as `axis side`; the lattice points of C as `boundary point` in the
%   accent colour, and p* (named pstar) as `marked point`;
% - the Dynkin diagram of AT Notation 3.7: a node at each lattice point of C'
%   not on C, joined along the boundary, black at lattice length 2 from p* and
%   white at length 1; at each corner node an internal node (`vertex=doubled`)
%   at the corner plus the primitive vectors of its two sides. Nodes are named
%   1, 2, ... along the boundary, then the internal nodes;
% - the monomial x^i y^j at every boundary lattice point but p*, and at every
%   internal node, set outside Q (along the outward normal of the side, or the
%   bisector of the two normals at a vertex) and, at an internal node, away
%   from its corner;
% - the decorated Dynkin symbol below Q.
% Keys: `ade labels=none` hides the monomials, side lengths and the p* label;
% `ade title=false` hides the symbol.

% x^i y^j, written 1, x, x^2, y, xy, ...
\newcommand{\dzgmonomial}[2]{%
  \ifnum#1=0 \ifnum#2=0 1\fi\else x\ifnum#1>1 ^{#1}\fi\fi
  \ifnum#2>0 y\ifnum#2>1 ^{#2}\fi\fi}

\newif\ifdzg@adelabels \dzg@adelabelstrue
\newif\ifdzg@adetitle \dzg@adetitletrue
\tikzset{
  ade labels/.is choice,
  ade labels/monomials/.code={\dzg@adelabelstrue},
  ade labels/none/.code={\dzg@adelabelsfalse},
  ade title/.is if=dzg@adetitle,
  pics/ade polygon/.style n args={3}{code={\dzg@adepolygon{#1}{#2}{#3}}},
}

% -----------------------------------------------------------------------------
% Shapes. Each sets p* = (\dzg@px,\dzg@py) and \dzg@path, the vertices
% v_1, ..., v_k of Q after p* in counterclockwise order: the boundary is
% p* -> v_1 -> ... -> v_k -> p*, with C = [v_k, p*] + [p*, v_1] and
% C' = v_1 -> ... -> v_k. For the affine shapes p* is interior to the one side
% of C. \dzg@lprime and \dzg@rprime record priming on the left side [p*, v_1]
% and the right side [v_k, p*].
% -----------------------------------------------------------------------------
\def\dzg@ade@rankerror#1{\PackageError{dzg}{ade polygon: #1}{AT Tables 1
  and 4 list the toric shapes.}}
\@namedef{dzg@ade@A@pure}{%
  \def\dzg@px{0}\def\dzg@py{2}%
  \pgfmathtruncatemacro\dzg@a{\dzg@r+1}\edef\dzg@path{0/0, \dzg@a/0}}
\@namedef{dzg@ade@A@left-short}{%
  \ifnum\dzg@r<1 \dzg@ade@rankerror{-A_n needs n >= 1}\fi
  \def\dzg@px{0}\def\dzg@py{2}%
  \pgfmathtruncatemacro\dzg@a{\dzg@r+2}\edef\dzg@path{1/0, \dzg@a/0}}
\@namedef{dzg@ade@A@prime}{%
  \ifnum\dzg@r<2 \dzg@ade@rankerror{'A_n needs n >= 2}\fi
  \def\dzg@lprime{1}%
  \pgfmathtruncatemacro\dzg@a{\dzg@r-1}\edef\dzg@path{0/1, 0/0, \dzg@a/0}}
\@namedef{dzg@ade@A@both-prime}{%
  \ifodd\dzg@r\else\dzg@ade@rankerror{'A'_n needs n odd}\fi
  \ifnum\dzg@r<5 \dzg@ade@rankerror{'A'_n needs n >= 5}\fi
  \def\dzg@lprime{1}\def\dzg@rprime{1}%
  \pgfmathtruncatemacro\dzg@a{\dzg@r-3}\pgfmathtruncatemacro\dzg@b{(\dzg@r+1)/2}%
  \edef\dzg@path{0/1, 0/0, \dzg@a/0, \dzg@b/1}}
\@namedef{dzg@ade@D@pure}{%
  \ifnum\dzg@r<4 \dzg@ade@rankerror{D_n needs n >= 4}\fi
  \pgfmathtruncatemacro\dzg@a{\dzg@r-2}\edef\dzg@path{0/2, 0/0, \dzg@a/0}}
\@namedef{dzg@ade@D@prime}{%
  \ifodd\dzg@r\dzg@ade@rankerror{D'_n needs n even}\fi
  \ifnum\dzg@r<6 \dzg@ade@rankerror{D'_n needs n >= 6}\fi
  \def\dzg@rprime{1}%
  \pgfmathtruncatemacro\dzg@a{\dzg@r-4}\pgfmathtruncatemacro\dzg@b{\dzg@r/2}%
  \edef\dzg@path{0/2, 0/0, \dzg@a/0, \dzg@b/1}}
\@namedef{dzg@ade@E@pure}{%
  \ifnum\dzg@r<6 \dzg@ade@rankerror{E_n needs 6 <= n <= 8}\fi
  \ifnum\dzg@r>8 \dzg@ade@rankerror{E_n needs 6 <= n <= 8}\fi
  \pgfmathtruncatemacro\dzg@a{\dzg@r-3}\edef\dzg@path{0/3, 0/0, \dzg@a/0}}
\@namedef{dzg@ade@D@affine}{%
  \ifodd\dzg@r\dzg@ade@rankerror{\tilde D_n needs n even}\fi
  \ifnum\dzg@r<6 \dzg@ade@rankerror{\tilde D_n is toric for n >= 6}\fi
  \def\dzg@affine{1}%
  \pgfmathtruncatemacro\dzg@a{\dzg@r-4}\edef\dzg@path{0/2, 0/0, \dzg@a/0, 4/2}}
\@namedef{dzg@ade@E@affine}{%
  \def\dzg@affine{1}%
  \ifcase\numexpr\dzg@r-6\relax
    \dzg@ade@rankerror{there is no \tilde E_6 shape (AT Remark 3.10)}%
  \or\def\dzg@path{0/4, 0/0, 4/0}%
  \or\def\dzg@path{0/3, 0/0, 6/0}%
  \else\dzg@ade@rankerror{\tilde E_n needs n = 7 or 8}\fi}

\def\dzg@ade@shape#1#2#3{%
  \def\dzg@letter{#1}\pgfmathtruncatemacro\dzg@r{#2}%
  \def\dzg@px{2}\def\dzg@py{2}\let\dzg@path\relax
  \def\dzg@lprime{0}\def\dzg@rprime{0}\def\dzg@affine{0}%
  \ifcsname dzg@ade@#1@#3\endcsname\csname dzg@ade@#1@#3\endcsname
  \else\dzg@ade@rankerror{no variant `#3' of #1}\fi}

% \dzg@ade@side{<x0>}{<y0>}{<x1>}{<y1>}: lattice length \dzg@g and primitive
% step (\dzg@sx,\dzg@sy) of the segment.
\def\dzg@ade@side#1#2#3#4{%
  \pgfmathtruncatemacro\dzg@g{gcd(abs(#3-#1),abs(#4-#2))}%
  \pgfmathtruncatemacro\dzg@sx{(#3-#1)/\dzg@g}%
  \pgfmathtruncatemacro\dzg@sy{(#4-#2)/\dzg@g}}

% \dzg@ade@monomial{<i>}{<j>}{<angle>}
\def\dzg@ade@monomial#1#2#3{%
  \ifdzg@adelabels
    \path (#1,#2) node[label={[monomial]#3:$\dzgmonomial{#1}{#2}$}] {};%
  \fi}

\def\dzg@adepolygon#1#2#3{%
  \dzg@ade@shape{#1}{#2}{#3}%
  \ifnum\dzg@affine=1 \let\dzg@cycle\dzg@path
  \else\edef\dzg@cycle{\dzg@px/\dzg@py, \dzg@path}\fi
  % Bounding box, the path of Q, and v_1, v_k.
  \gdef\dzg@w{0}\gdef\dzg@h{0}\gdef\dzg@Q{}%
  \foreach \x/\y [count=\k] in \dzg@cycle {%
    \xdef\dzg@N{\k}%
    \pgfmathtruncatemacro\dzg@t{max(\dzg@w,\x)}\xdef\dzg@w{\dzg@t}%
    \pgfmathtruncatemacro\dzg@t{max(\dzg@h,\y)}\xdef\dzg@h{\dzg@t}%
    \xdef\dzg@Q{\dzg@Q(\x,\y) -- }}%
  \foreach \x/\y [count=\k] in \dzg@path {%
    \ifnum\k=1 \xdef\dzg@vfirst{\x/\y}\xdef\dzg@vfirstx{\x}\xdef\dzg@vfirsty{\y}\fi
    \xdef\dzg@vlast{\x/\y}\xdef\dzg@vlastx{\x}\xdef\dzg@vlasty{\y}}%
  \path[local bounding box=lp] (0,0) rectangle (\dzg@w,\dzg@h);
  \begin{scope}[on background layer]
    \fill[polygon region] \dzg@Q cycle;
    \draw[lattice grid] (0,0) grid (\dzg@w,\dzg@h);
  \end{scope}
  \foreach \x in {0,...,\dzg@w} \foreach \y in {0,...,\dzg@h}
    \node[lattice point] (\x-\y) at (\x,\y) {};
  % Monomials: outward along the normal of each side at its interior points,
  % along the bisector of the two outward normals at each vertex.
  \edef\dzg@wrapped{\dzg@cycle, \dzg@cycle}%
  \foreach \x/\y [count=\k, remember=\xa as \xb (initially 0),
      remember=\ya as \yb (initially 0), remember=\x as \xa (initially 0),
      remember=\y as \ya (initially 0)] in \dzg@wrapped {%
    \pgfmathtruncatemacro\dzg@atvertex{\k>2 && \k<\dzg@N+3}%
    \pgfmathtruncatemacro\dzg@atside{\k>1 && \k<\dzg@N+2}%
    \ifnum\dzg@atvertex=1 % vertex (\xa,\ya) between (\xb,\yb) and (\x,\y)
      \pgfmathsetmacro\dzg@nx{(\ya-\yb)/veclen(\xa-\xb,\ya-\yb)
        + (\y-\ya)/veclen(\x-\xa,\y-\ya)}%
      \pgfmathsetmacro\dzg@ny{(\xb-\xa)/veclen(\xa-\xb,\ya-\yb)
        + (\xa-\x)/veclen(\x-\xa,\y-\ya)}%
      \pgfmathsetmacro\dzg@angle{atan2(\dzg@ny,\dzg@nx)}%
      \ifnum\xa=\dzg@px \ifnum\ya=\dzg@py \else
        \dzg@ade@monomial{\xa}{\ya}{\dzg@angle}\fi
      \else\dzg@ade@monomial{\xa}{\ya}{\dzg@angle}\fi
    \fi
    \ifnum\dzg@atside=1 % interior lattice points of side (\xa,\ya) -> (\x,\y)
      \dzg@ade@side{\xa}{\ya}{\x}{\y}%
      \pgfmathsetmacro\dzg@angle{atan2(-\dzg@sx,\dzg@sy)}%
      \pgfmathtruncatemacro\dzg@last{\dzg@g-1}%
      \ifnum\dzg@last>0
        \foreach \t in {1,...,\dzg@last} {%
          \pgfmathtruncatemacro\dzg@i{\xa+\t*\dzg@sx}%
          \pgfmathtruncatemacro\dzg@j{\ya+\t*\dzg@sy}%
          \ifnum\dzg@i=\dzg@px \ifnum\dzg@j=\dzg@py \else
            \dzg@ade@monomial{\dzg@i}{\dzg@j}{\dzg@angle}\fi
          \else\dzg@ade@monomial{\dzg@i}{\dzg@j}{\dzg@angle}\fi}%
      \fi
    \fi}%
  % C: the sides through p*, and their lattice points.
  \ifnum\dzg@affine=1
    \edef\dzg@Csides{\dzg@vlast/\dzg@vfirst}%
  \else
    \edef\dzg@Csides{\dzg@vlast/\dzg@px/\dzg@py, \dzg@px/\dzg@py/\dzg@vfirst}%
  \fi
  \foreach \xa/\ya/\x/\y in \dzg@Csides {%
    \dzg@ade@side{\xa}{\ya}{\x}{\y}%
    \ifodd\dzg@g\def\dzg@class{short side}\else\def\dzg@class{long side}\fi
    \scoped[on edge layer] \draw[\dzg@class] (\xa,\ya) -- (\x,\y);
    \ifdzg@adelabels \path[side label=\dzg@g] (\xa,\ya) -- (\x,\y);\fi
    \foreach \t in {0,...,\dzg@g}
      \node[boundary point=dzg accent] at (\xa+\t*\dzg@sx,\ya+\t*\dzg@sy) {};}%
  \coordinate (pstar) at (\dzg@px,\dzg@py);
  \ifdzg@adelabels \node[marked point=$p_\ast$] at (pstar) {};
  \else \node[marked point] at (pstar) {};\fi
  % C' and the Dynkin diagram along it.
  \gdef\dzg@node{0}%
  \foreach \x/\y [count=\k, remember=\x as \xa (initially 0),
      remember=\y as \ya (initially 0)] in \dzg@path {%
    \ifnum\k>1
      \scoped[on edge layer] \draw[axis side] (\xa,\ya) -- (\x,\y);
      \dzg@ade@side{\xa}{\ya}{\x}{\y}%
      \foreach \t in {1,...,\dzg@g} {%
        \pgfmathtruncatemacro\dzg@i{\xa+\t*\dzg@sx}%
        \pgfmathtruncatemacro\dzg@j{\ya+\t*\dzg@sy}%
        \edef\dzg@here{\dzg@i/\dzg@j}%
        \ifx\dzg@here\dzg@vlast\else
          \pgfmathtruncatemacro\dzg@n{\dzg@node+1}\xdef\dzg@node{\dzg@n}%
          \pgfmathtruncatemacro\dzg@len{gcd(abs(\dzg@i-\dzg@px),abs(\dzg@j-\dzg@py))}%
          \ifnum\dzg@len=2 \def\dzg@mark{black}\else\def\dzg@mark{white}\fi
          \node[vertex=\dzg@mark] (\dzg@n) at (\dzg@i,\dzg@j) {};
          \ifnum\dzg@n>1
            \scoped[on edge layer]
              \draw[coxeter edge=3] (\dzg@i-\dzg@sx,\dzg@j-\dzg@sy) -- (\dzg@i,\dzg@j);
          \fi
        \fi}%
    \fi}%
  % Internal nodes at the corners of C'.
  \foreach \x/\y [count=\k, remember=\xa as \xb (initially 0),
      remember=\ya as \yb (initially 0), remember=\x as \xa (initially 0),
      remember=\y as \ya (initially 0)] in \dzg@path {%
    \ifnum\k>2 % corner (\xa,\ya)
      \dzg@ade@side{\xa}{\ya}{\xb}{\yb}\let\dzg@ux\dzg@sx\let\dzg@uy\dzg@sy
      \dzg@ade@side{\xa}{\ya}{\x}{\y}%
      \pgfmathtruncatemacro\dzg@dx{\dzg@ux+\dzg@sx}%
      \pgfmathtruncatemacro\dzg@dy{\dzg@uy+\dzg@sy}%
      \pgfmathtruncatemacro\dzg@i{\xa+\dzg@dx}%
      \pgfmathtruncatemacro\dzg@j{\ya+\dzg@dy}%
      \pgfmathtruncatemacro\dzg@n{\dzg@node+1}\xdef\dzg@node{\dzg@n}%
      \node[vertex=doubled] (\dzg@n) at (\dzg@i,\dzg@j) {};
      \scoped[on edge layer] \draw[coxeter edge=3] (\xa,\ya) -- (\dzg@i,\dzg@j);
      \pgfmathsetmacro\dzg@angle{atan2(\dzg@dy,\dzg@dx)}%
      \dzg@ade@monomial{\dzg@i}{\dzg@j}{\dzg@angle}%
    \fi}%
  % The decorated Dynkin symbol (AT Notations 3.8, 3.21): a prime on a primed
  % side, else a minus on a short side; an affine shape carries its one
  % decoration on the right.
  \ifdzg@adetitle
    \dzg@ade@side{\dzg@px}{\dzg@py}{\dzg@vfirstx}{\dzg@vfirsty}%
    \let\dzg@lg\dzg@g
    \dzg@ade@side{\dzg@vlastx}{\dzg@vlasty}{\dzg@px}{\dzg@py}%
    \let\dzg@rg\dzg@g
    \ifnum\dzg@affine=1
      \dzg@ade@side{\dzg@vlastx}{\dzg@vlasty}{\dzg@vfirstx}{\dzg@vfirsty}%
      \let\dzg@rg\dzg@g\def\dzg@lg{2}%
      \def\dzg@body{\widetilde{\dzg@letter}}%
    \else\def\dzg@body{\dzg@letter}\fi
    \ifnum\dzg@lprime=1 \def\dzg@pre{{}'}%
    \else\ifodd\dzg@lg\def\dzg@pre{{}^{-}}\else\def\dzg@pre{}\fi\fi
    \ifnum\dzg@rprime=1 \def\dzg@post{'}%
    \else\ifodd\dzg@rg\def\dzg@post{^{-}}\else\def\dzg@post{}\fi\fi
    \node[polygon title, below=7mm of lp.south]
      {$\dzg@pre\dzg@body_{\dzg@r}\dzg@post$};
  \fi}

% =============================================================================
% Complete fans of toric varieties
% =============================================================================
%   \pic (S) {toric fan=<name>};
% draws the fan's rays from the origin S-0 to S-rho<i> (the primitive
% generator times `fan length`, default 2), labels each ray rho_i = (a,b),
% and names the maximal cone sigma_i, which is labelled at S-sigma<i>.
% `fan labels=none` hides the labels. Figures add lattice points, polytopes
% and shading around the named points.
%   P1   the fan of P^1 in N_R = R: sigma_0 = R_{>=0} = Cone(rho_0),
%        sigma_1 = R_{<=0} = Cone(rho_1), the normal fan of an interval
%        [CLS Example 2.4.10]
%   F2   the fan of the Hirzebruch surface F_2, the normal fan of the lattice
%        quadrilateral P_{2,4} = conv{(0,0), (2,0), (4,1), (0,1)}
%        [CLS Example 2.3.16]: rays rho_1..rho_4 = (1,0), (0,1),
%        (-1,2), (0,-1) counterclockwise, sigma_i = Cone(rho_i, rho_{i+1})
\newif\ifdzg@fanlabels \dzg@fanlabelstrue
\def\dzg@fanlength{2}
\tikzset{
  fan labels/.is choice,
  fan labels/canonical/.code={\dzg@fanlabelstrue},
  fan labels/none/.code={\dzg@fanlabelsfalse},
  fan length/.store in=\dzg@fanlength,
  pics/toric fan/.style={code={%
    \ifcsname dzg@fan@#1\endcsname\else
      \PackageError{dzg}{toric fan=#1 is not defined}{The fans are P1, F2.}\fi
    \node[distinguished point] (-0) at (0,0) {};
    \csname dzg@fan@#1\endcsname}},
}
% \dzg@fanrays{<index>/<a>/<b>/<label anchor>, ...}
\def\dzg@fanrays#1{%
  \foreach \dzg@i/\dzg@a/\dzg@b/\dzg@p in {#1}{%
    \coordinate (-rho\dzg@i) at (\dzg@fanlength*\dzg@a, \dzg@fanlength*\dzg@b);
    \ifdzg@fanlabels
      \draw[fan ray] (-0) -- (-rho\dzg@i)
        node[\dzg@p, font=\small] {$\rho_{\dzg@i} = (\dzg@a,\dzg@b)$};
    \else
      \draw[fan ray] (-0) -- (-rho\dzg@i);
    \fi}}
% \dzg@fancones{<index>/<coordinate>, ...}
\def\dzg@fancones#1{%
  \foreach \dzg@i/\dzg@c in {#1}{%
    \coordinate (-sigma\dzg@i) at (\dzg@c);
    \ifdzg@fanlabels \node[font=\small] at (-sigma\dzg@i) {$\sigma_{\dzg@i}$};\fi}}
\@namedef{dzg@fan@P1}{%
  \foreach \dzg@i/\dzg@s/\dzg@p in {0/1/right, 1/-1/left}{%
    \coordinate (-rho\dzg@i) at (\dzg@s*\dzg@fanlength, 0);
    \draw[fan ray] (-0) -- (-rho\dzg@i);
    \coordinate (-sigma\dzg@i) at (\dzg@s*0.6*\dzg@fanlength, 0);
    \ifdzg@fanlabels
      \node[below, font=\small] at (-sigma\dzg@i) {$\sigma_{\dzg@i}$};\fi}}
\@namedef{dzg@fan@F2}{%
  \dzg@fanrays{1/1/0/right, 2/0/1/above right, 3/-1/2/above left, 4/0/-1/below}%
  \dzg@fancones{1/{0.35*\dzg@fanlength,0.35*\dzg@fanlength},
    2/{-0.25*\dzg@fanlength,0.7*\dzg@fanlength},
    3/{-0.4*\dzg@fanlength,-0.3*\dzg@fanlength},
    4/{0.35*\dzg@fanlength,-0.3*\dzg@fanlength}}}
